% Generated by code/rebuild/WriteResults.py from corrected-run outputs.
\section{Results}

We first asked whether a lasting change in how some traders respond to
their neighbors could push a market of mixed Buy, Hold, and Sell
decisions to a Sell majority.
Rule~19286 differs from the baseline Rule~19280 in a single response:
a trader with three neighbors in Buy and one in Hold chooses Sell rather
than Buy (Fig.~\ref{fig:rule-quench}A,B). We paired each switched market
with a no-switch control that kept Rule~19280 and shared its baseline
history (Supplementary
Table~\ref{tab:campaigns}). With active technical traders and dealer
depth $h=100$, none of the 50 starting markets reached a Sell majority
within 3,000 steps after the switch when the fraction of interior
traders given the new rule was $f=0.10$ or 0.15. All of them did at
$f=0.20$ and 0.25, whereas no control reached a majority over the same
period. The median onset fell from 1,070 steps after the switch at
fraction $f=0.20$ to 311.5 steps at $f=0.25$
(Supplementary Table~\ref{tab:paired-switch}). In the displayed pair at
fraction $f=0.20$, groups of sellers expanded after the rule switch
while the control kept its mixture of decisions
(Fig.~\ref{fig:rule-quench}C). To test whether a single burst of Sell
intentions could do the same, we instead set the intentions of the same
selected traders to Sell once at the switch and let them keep
Rule~19280 for every later update. This one-time
instruction produced no majority at fraction $f=0.15$, 0.20, or 0.25,
whether or not the technical traders responded to price
(Supplementary Table~\ref{tab:revision-pulse}). Which response changed
also mattered. We screened all eighteen rules that differ from
Rule~19280 in a single response, including Rule~19286, over 20 starting
markets with technical traders active and held. At fraction $f=0.20$,
six rules produced a majority in every market under both conditions.
Five of them placed Sell at one of the neighbor sums 1 to 5, the only sums
at which Rule~19280 does not already assign Sell, and Rule~19286 was one
of these five. The sixth replaced Buy with Hold at neighbor sum 3. The other twelve produced no majority in any market: ten moved a response away from Sell, and two replaced Buy with Hold at
neighbor sums 1 and 2 (Fig.~\ref{fig:rule-quench}B;
Supplementary Table~\ref{tab:revision-rules}). Thus, the lasting
switch to Rule~19286 produced majorities where a one-time Sell
instruction to the same traders did not, and outcomes depended on which
response changed and how many traders received it.

The rule screen produced majorities with technical traders held as
well as active, which suggested that a Sell majority did not require
price feedback. To test this across
fractions and dealer depths, we repeated the rule switch at 11 fractions
between 0.10 and 0.25 with the technical traders either active or held
in Hold (Fig.~\ref{fig:reservoir}). We also varied the dealer depth over
$h=0$, 0.1, 0.5, 1, 2, 5, 10, and 100, because depth sets how far the
quoted price moves for a given net trade. In the held condition,
trades still completed and, at positive depths, the price still moved,
but no decision responded to it. Because no decision in the held
condition depended on cash, shares, or price, the onset, final mean Sell
fraction, and final interior grid for each starting market and fraction
were identical across dealer depths despite different price paths. The
held Sell-majority probability was therefore the same at every depth:
0.34 at fraction $f=0.175$, 0.70 at $f=0.185$, and
1.00 at $f=0.20$ (Fig.~\ref{fig:reservoir}A). Responses to neighbors
were thus sufficient to produce a Sell majority.
Majority probability rose from zero to one over a narrow range of
fractions, which we refer to as the crossover. Within it, the
probability was usually higher with technical traders active than
held. At dealer depth $h=100$, the probability
with technical traders active was 0.50 at fraction $f=0.175$, compared
with 0.34 when held, and 0.84 at $f=0.185$, compared with 0.70
(Fig.~\ref{fig:reservoir}A). Across the 88 combinations of eight depths and 11 fractions, the active
minus held difference was positive in 22, with a largest value of 0.18,
and negative in two, by at most 0.04. It was zero in the rest, including
every fraction outside the crossover and every combination at $h=0$,
where the price never moved, the technical traders' band kept zero
width, and they stayed in Hold
(Fig.~\ref{fig:reservoir}B). Thus, price feedback could change how
often the rule switch produced a majority, but it was not needed for
selling decisions to spread.

\input{sections/finite-size-results}

Because the midpoint moved with observation time, we extended five
example markets with active technical traders and dealer depth $h=100$
to 10,000 steps after the switch, to see whether markets without a
majority at 3,000 steps reached one later and how the others developed
(Fig.~\ref{fig:dose-panels}). Before any grid
images were inspected, we had selected them from the 3,000-step outcomes
as typical of their fraction and outcome. They comprised one market
without a majority at fraction $f=0.15$, one with and one without a
majority at $f=0.175$, and one with a majority at each of $f=0.185$ and
0.20. The market without a majority at fraction $f=0.15$ still had none
after 10,000 steps, whereas the market without one at $f=0.175$ reached
a majority 4{,}073 steps after the switch. That market's maximum
drawdown rose
from 10.17\% at 3,000 steps to 26.34\% at 10,000 steps
(Fig.~\ref{fig:dose-panels}B; Supplementary
Table~\ref{tab:market-panel-horizons}). The other three majorities had formed
2{,}330, 1{,}836, and 1{,}028 steps after the switch at fractions
$f=0.175$, 0.185, and 0.20, respectively (Supplementary
Table~\ref{tab:market-panel-horizons}). Each displayed majority formed
as groups of sellers grew and merged (Fig.~\ref{fig:dose-panels}A,C).
At fraction $f=0.15$, no switched market had reached a majority by
3,000 steps, so any majority there would form late. To measure how
often this happened, we extended to 10,000 steps the runs from all 50
starting markets at that fraction, with the same active technical
traders and dealer depth, together with their no-switch controls. Two
switched markets reached a majority by 10,000 steps, compared with none
at 3,000 steps, and no control reached one. The longer runs also showed
how the switch affected prices before any majority formed. Every
control stayed below a maximum drawdown of 1.15\% over 10,000 steps.
Among the 50 switched markets at this fraction, the median drawdown was
already 5.04\% at 3,000 steps, before any of them had a majority, and
5.87\% at 10,000 steps. Their largest drawdown was 12.29\% at 3,000
steps. At 10,000 steps it was 25.30\%, in a different market, one that first reached a majority between 3,000
and 10,000 steps
(Supplementary Table~\ref{tab:market-panel-horizons}). Thus, remaining
without a majority for 3,000 steps did not establish a long-run state,
the rule switch deepened price declines even before any majority formed,
and a majority that formed later could raise a market's drawdown well
above its value at 3,000 steps.

We next asked how dealer reserves affected the price decline that
accompanied a majority, because a dealer with small reserves might be
expected to limit a decline by buying fewer of the shares that traders
offered. Using the same dealer-depth comparison, we examined the quoted
price after the rule switch (Fig.~\ref{fig:reservoir}). At fraction
$f=0.20$, every starting market reached a majority at every tested
depth, with technical traders active or held. The declines
nonetheless differed in size: with active technical traders, the median maximum drawdown of switched
markets was 99.99\% at dealer depth $h=0.1$, 99.52\% at $h=1$, and
25.15\% at $h=100$ (Fig.~\ref{fig:reservoir}C). Additional runs
at dealer depths $h=40$ and 400, outside the eight-depth sweep, at the
same fraction and with active technical traders, followed the same
decrease with depth, with median drawdowns of 48.62\% and 7.08\%
(Supplementary Table~\ref{tab:revision-impact}). Smaller
reserves made the quoted price more sensitive to net selling, because
the same net sale added a larger proportion to a small share reserve.
Shallow positive reserves did not limit these declines by refusing
trades. A constant-product dealer never ran out of cash for buying, and
its few rejections came from clearing limits: median dealer-rejection
fractions were 226, 34, and 0.6 per million eligible orders at dealer
depths $h=0.1$, 1, and 100 (Fig.~\ref{fig:reservoir}D).
The same price sensitivity also produced large declines without the
rule switch.
At shallow depth, ordinary order imbalances moved the quoted price up and
down by large amounts, so the largest fall from an earlier peak was large
even without a lasting trend.
No-switch controls at dealer depth $h=0.1$ had a median drawdown of
77.04\%, compared with 0.75\% at $h=100$. The median over starting markets of each switched market's increase in
drawdown over its no-switch control rose from 22.96 percentage points at
dealer depth $h=0.1$ to a largest value of 80.31 percentage points at
$h=5$ and fell to 24.39 percentage points at $h=100$ (Supplementary
Table~\ref{tab:paired-switch}). At $h=5$, the median drawdown was
14.03\% in controls and 94.43\% in switched markets.
The increase was smaller at both ends because at $h=0.1$ the controls'
decline left little room for more, whereas at $h=100$ the quoted price
responded weakly to net selling. Zero reserves were a separate case, in which only opposing trader orders
matched, the dealer did not trade, and the price stayed fixed, so the
drawdown was zero although every market reached a majority. In the
displayed starting market 2 at fraction $f=0.20$, a majority formed at
dealer depths $h=0.1$, 1, and 100 with active technical traders, while
the quoted price fell nearly to zero at the two shallower depths and by
about a quarter at $h=100$ (Fig.~\ref{fig:dealer-panels}). Thus, at
fraction $f=0.20$, the dealer's reserves set how far the quoted price
fell both with and without the rule switch, and shallow positive
reserves did not prevent severe drawdowns.

Trades also depended on the traders' own resources, so we independently
multiplied their cash and shares by 0.1, 1, or 10 at the rule switch,
at dealer depths $h=0.1$, 1, and 100 (Fig.~\ref{fig:balance-sheet}). We
report fraction $f=0.18$, near the held midpoint of 0.178 for the
$64\times64$ grid used here, and give three neighboring fractions in Supplementary
Table~\ref{tab:balance-fractions}. Holdings entered the order and
clearing calculations but not the interior decision rule, so they could
change decisions only through the price response of active technical
traders. With technical traders held, every cash and share combination
produced the same decision history within a starting market, and the
Sell-majority probability was 0.64 at every dealer depth, with a 95\%
Wilson interval of [0.50, 0.76]. With technical traders active, the
probability was 0.58 to 0.60 with unchanged resources, slightly below
the held value, and ranged from 0.54 to 0.80 across all resource
combinations and dealer depths (Fig.~\ref{fig:balance-sheet}A). This comparison had no separate
no-switch controls, so its drawdowns describe switched markets rather
than effects of the switch. Among the 27 active combinations, the five
highest majority probabilities all had the share multiplier 0.1, and
within each combination of cash and dealer depth that multiplier also
gave the smallest drawdown, so a higher majority probability did not
bring a larger drawdown (Fig.~\ref{fig:balance-sheet}D). With unchanged cash
and dealer depth $h=100$, share multipliers of 0.1, 1, and 10 gave
median drawdowns of 3.17\%, 20.18\%, and 77.55\%
(Fig.~\ref{fig:balance-sheet}B). Drawdowns rose with shares because
fewer Sell orders were blocked for lack of shares, so more net selling
reached the dealer's reserves, which the intervention did not rescale. By contrast, raising the cash
multiplier from 1 to 10 at dealer depth $h=1$ with unchanged shares
changed the median drawdown only from 96.46\% to 97.10\% and the
majority probability from 0.58 to 0.56, in markets where the decline was
already close to complete (Fig.~\ref{fig:balance-sheet}A,B). Dealer
rejection did not cushion these declines: median dealer-rejection
fractions stayed below 5,000 per million eligible orders in every
combination, including those with the largest drawdowns
(Fig.~\ref{fig:balance-sheet}B,C). Together, the dealer-depth and
resource comparisons showed that financial outcomes in switched markets
depended on traders' cash and shares and on the dealer's reserves, so a
Sell majority did not by itself determine how far the quoted price
would fall.
